El sistema se utiliza desde una computadora conectada al puerto USB--UART de la
Basys 3. Desde allí se envían los operandos y el opcode de cada operación. La
placa procesa esos valores y devuelve el resultado por el mismo enlace.

Los datos no pasan directamente desde la computadora hasta la ALU. Primero se
transmiten por UART como una secuencia de tres bytes. El controlador recupera
esos valores y los presenta en paralelo a la ALU. Luego, el resultado vuelve a
la computadora y también queda visible en los LED.

La integración conecta los bloques que participan en ese recorrido. En este
nivel se establece cómo se transfieren los datos, cuándo se acepta cada byte y
qué información se presenta fuera de la FPGA. Cada módulo mantiene una función
propia y se comunica con los demás mediante señales bien definidas.

\section{Recorrido de una transacción}

Cada solicitud comienza con el envío del operando A, el operando B y el opcode.
\texttt{uart\_core} recibe las tres tramas y almacena los bytes reconstruidos.
El controlador los retira en el mismo orden y guarda cada uno en el registro
correspondiente.

Cuando la solicitud está completa, la ALU recibe ambos operandos y los seis bits
inferiores del opcode. El controlador registra el resultado y las banderas en
el ciclo siguiente. El resultado se entrega entonces a la cola de transmisión
y los LED se actualizan con la información de la operación.

La tabla~\ref{tab:secuencia-operacion} resume este recorrido. Cada transferencia
entre bloques se realiza en un flanco del reloj.

\begin{table}[H]
	\centering
	\begin{tabularx}{0.96\textwidth}{C{1.1cm}L{4.2cm}Y}
		\toprule
		\textbf{Paso} & \textbf{Transferencia} & \textbf{Efecto} \\
		\midrule
		0 & Reset hacia todos los bloques &
		Las colas quedan vacías y el resultado inicial toma el valor cero. \\
		1 & Computadora hacia \texttt{uart\_core} &
		Las tramas de A, B y opcode se reconstruyen como bytes. \\
		2 & FIFO RX hacia el controlador &
		Los tres campos se guardan en el orden definido por el protocolo. \\
		3 & Controlador hacia la ALU &
		A, B y opcode quedan disponibles en paralelo. \\
		4 & ALU hacia el controlador &
		El resultado y las banderas se registran y actualizan los LED. \\
		5 & Controlador hacia FIFO TX &
		El byte de resultado queda pendiente de transmisión. \\
		6 & \texttt{uart\_tx} hacia la computadora &
		El resultado se serializa y se envía por la línea UART. \\
		\bottomrule
	\end{tabularx}
	\caption{Recorrido de una operación a través del sistema integrado.}%
	\label{tab:secuencia-operacion}
\end{table}

El controlador vuelve a esperar una solicitud después de registrar la operación
anterior. De esta manera, el byte transmitido y el patrón de los LED
corresponden
a los mismos operandos y al mismo opcode.

\section{Resultado UART e indicadores}

La respuesta UART contiene solamente los ocho bits del resultado. Las banderas
Z, C y V no se transmiten. En cambio, se muestran junto con el resultado en los
once LED de la placa.

El controlador registra todos estos valores al completar una operación. La
salida hacia los LED se forma a partir de esos registros según la siguiente
distribución:

\begin{equation}
	\texttt{leds[10:0]} =
	\{\mathrm{V},\,\mathrm{C},\,\mathrm{Z},\,\mathrm{resultado[7:0]}\}.
\end{equation}

Los valores registrados se mantienen hasta que termina una nueva operación.
Esto permite observar el resultado aunque la respuesta UART ya haya sido
transmitida. Después del reset, los ocho bits de resultado quedan en cero. La
bandera Z queda activa y las banderas C y V permanecen apagadas.

\section{Coordinación y rendimiento}

Los bloques no siempre producen y consumen los bytes en el mismo ciclo. La cola
de recepción conserva los datos que todavía no fueron leídos por el controlador.
La cola de transmisión cumple la misma función con los resultados que todavía
no fueron enviados.

El controlador solo retira un byte cuando la cola de recepción contiene datos.
También comprueba que exista espacio antes de guardar una respuesta. Si la cola
de transmisión está llena, conserva el resultado y espera hasta que pueda
escribirlo. Así, una diferencia momentánea entre los bloques no provoca la
pérdida de información.

La velocidad máxima queda determinada por el enlace UART. Cada trama 8N1 ocupa
diez bits. A \num{19200} baud pueden transmitirse \num{1920} bytes por segundo.
Como cada solicitud contiene tres bytes, la entrada admite como máximo
\num{640} operaciones por segundo.

Cada operación genera un solo byte de respuesta. Por lo tanto, esa cantidad de
solicitudes requiere \num{6400} bits por segundo en el sentido de transmisión.
Este valor se encuentra por debajo de la capacidad disponible. Además, las dos
líneas UART permiten recibir una solicitud mientras se transmite el resultado
anterior.

El cliente espera la respuesta de una operación antes de enviar la siguiente.
Con esta forma de intercambio, cada resultado puede asociarse directamente con
la solicitud que lo produjo.

\section{Cliente de la computadora}

Para enviar las operaciones desde la computadora se desarrolló un programa que
actúa como cliente del sistema. Este programa evita que el usuario tenga que
configurar el puerto y transmitir cada byte de forma manual.

El protocolo trabaja con valores binarios y no con caracteres de texto. El
cliente \texttt{tools/uart\allowbreak\_alu\allowbreak\_client.py} configura el puerto serie en 8N1 a
\num{19200} baud. Luego forma la solicitud con A, B y opcode, y envía los tres
bytes en ese orden.

Después del envío, el cliente recibe un único byte y lo compara con un modelo de
referencia. Puede ejecutar una operación individual o una serie de casos
dirigidos y pseudoaleatorios. El modelo también calcula las banderas para poder
compararlas con los LED de la placa.

El cliente incluye además una prueba de recuperación. Para realizarla, envía un
byte aislado y espera \SI{20}{\milli\second}. Este intervalo supera el timeout
del controlador y hace que la solicitud incompleta sea descartada. A
continuación, se envía una operación completa y se comprueba que la respuesta
corresponda a esa nueva solicitud.

\section{Integración con la Basys 3}

El módulo superior \texttt{alu\_uart\_top.sv} conecta el diseño con
los recursos de la placa. Sus puertos corresponden al reloj de
entrada, el pulsador de reset, las dos líneas UART y los once LED. El
archivo \texttt{basys3.xdc} de restricciones asigna estas señales a
los pines de la Basys 3.

\begin{table}[H]
	\centering
	\begin{tabularx}{0.92\textwidth}{L{3.5cm}C{2.0cm}Y}
		\toprule
		\textbf{Puerto} & \textbf{Pin} & \textbf{Recurso físico} \\
		\midrule
		\texttt{clk} & W5 & Oscilador de \SI{100}{\mega\hertz}. \\
		\texttt{btn\_reset} & U18 & Pulsador central BTNC. \\
		\texttt{serial\_rx} & B18 & USB--UART desde la computadora. \\
		\texttt{serial\_tx} & A18 & USB--UART hacia la computadora. \\
		\texttt{leds[7:0]} & varios & Resultado de ocho bits. \\
		\texttt{leds[10:8]} & W3, V3, V13 & Banderas V, C y Z. \\
		\bottomrule
	\end{tabularx}
	\caption{Correspondencia entre los puertos del diseño y la Basys 3.}%
\end{table}

Las señales \texttt{serial\_rx} y \texttt{btn\_reset} se originan fuera de la
FPGA y pueden cambiar en cualquier momento respecto del reloj interno. Por este
motivo, ambas reciben un tratamiento especial antes de utilizarse en el resto
del circuito.

La entrada \texttt{serial\_rx} atraviesa el sincronizador ubicado antes del
receptor UART. En el caso del reset, la activación es inmediata para llevar los
bloques a un estado conocido. Su liberación se sincroniza con el reloj para que
todos retomen el funcionamiento a partir de un flanco válido.

Estas características también se tuvieron en cuenta al definir las restricciones
temporales. El archivo \texttt{basys3.xdc} establece un período de
\SI{10}{\nano\second} para el reloj y aplica las excepciones únicamente al
recorrido de las entradas asíncronas hasta sus sincronizadores. A partir de ese
punto, las señales se analizan del mismo modo que el resto de la lógica interna.
El archivo también establece el estándar eléctrico LVCMOS33 para los puertos de
la placa.

Con estas conexiones, el resultado de una operación puede comprobarse de dos
formas. La computadora recibe el byte transmitido por UART y la placa muestra
ese mismo valor junto con las banderas.
